\documentclass{article}

\usepackage[utf8]{inputenc}
\usepackage[T2A]{fontenc}
\usepackage[russian]{babel}

\begin{document}
\small

\section*{Task 1.3}

\section*{Возможные трассы}

\subsection*{1) $(x=1,y=0,z=1)$}

\[
A.1 \to A.2 \to A.3 \to B.1 \to B.2 \to B.3 \to B.4 \to C.1 \to C.2
\]

\subsection*{2) $(x=1,y=1,z=1)$}

\[
B.1 \to B.2 \to A.1 \to A.2 \to A.3 \to B.3 \to B.4 \to C.1 \to C.2
\]

\subsection*{3) $(x=1,y=2,z=1)$}

\[
B.1 \to B.2 \to B.3 \to B.4 \to C.1 \to C.2 \to A.1 \to A.2 \to A.3
\]

\subsection*{4) $(x=0,y=2,z=1)$}

\[
B.1 \to B.2 \to B.3 \to B.4 \to A.1 \to A.2 \to A.3 \to C.1 \to C.2 \to C.3 \to C.4
\]

\section*{Доказательство, что других трасс не может быть}

\subsection*{Фаза 1}

\begin{itemize}
\item $z$ может изменится только в $B.4$, поэтому всегда $z=1$.
\item $x$ меняется в $B.2$ и в $C.4$ (на +1/-1 соответственно) и максимум 1 раз. Значит $x \in \{0,1\}$.
\item $y=a_x+a_z$, где $a_x,a_z \in \{0,1\}$. Значит $y \in \{0,1,2\}$.
\end{itemize}

Таким образом, получили границы, где переменные принадлежат множествам: $z=1$,
$x \in \{0,1\}$,
$y \in \{0,1,2\}$.

\subsection*{Фаза 2}

\begin{itemize}
\item $y=0$ возможно только если A читает старые $x=0,z=0$.
\item $y=1$ если один из них уже стал 1.
\item $y=2$ если A читает $x=1,z=1$.
\item $x=0$ возможно только если было $C.4$, значит в C было $y=2$.
\item Если $C$ читает $y\neq2$, значит уменьшения нет и $x=1$.
\end{itemize}

Других комбинаций быть не может.

\end{document}